% Options for packages loaded elsewhere
\PassOptionsToPackage{unicode}{hyperref}
\PassOptionsToPackage{hyphens}{url}
\documentclass[
  10pt,
]{article}
\usepackage{xcolor}
\usepackage[margin=2cm]{geometry}
\usepackage{amsmath,amssymb}
\setcounter{secnumdepth}{5}
\usepackage{iftex}
\ifPDFTeX
  \usepackage[T1]{fontenc}
  \usepackage[utf8]{inputenc}
  \usepackage{textcomp} % provide euro and other symbols
\else % if luatex or xetex
  \usepackage{unicode-math} % this also loads fontspec
  \defaultfontfeatures{Scale=MatchLowercase}
  \defaultfontfeatures[\rmfamily]{Ligatures=TeX,Scale=1}
\fi
\usepackage{lmodern}
\ifPDFTeX\else
  % xetex/luatex font selection
\fi
% Use upquote if available, for straight quotes in verbatim environments
\IfFileExists{upquote.sty}{\usepackage{upquote}}{}
\IfFileExists{microtype.sty}{% use microtype if available
  \usepackage[]{microtype}
  \UseMicrotypeSet[protrusion]{basicmath} % disable protrusion for tt fonts
}{}
\makeatletter
\@ifundefined{KOMAClassName}{% if non-KOMA class
  \IfFileExists{parskip.sty}{%
    \usepackage{parskip}
  }{% else
    \setlength{\parindent}{0pt}
    \setlength{\parskip}{6pt plus 2pt minus 1pt}}
}{% if KOMA class
  \KOMAoptions{parskip=half}}
\makeatother
\usepackage{color}
\usepackage{fancyvrb}
\newcommand{\VerbBar}{|}
\newcommand{\VERB}{\Verb[commandchars=\\\{\}]}
\DefineVerbatimEnvironment{Highlighting}{Verbatim}{commandchars=\\\{\}}
% Add ',fontsize=\small' for more characters per line
\usepackage{framed}
\definecolor{shadecolor}{RGB}{248,248,248}
\newenvironment{Shaded}{\begin{snugshade}}{\end{snugshade}}
\newcommand{\AlertTok}[1]{\textcolor[rgb]{0.94,0.16,0.16}{#1}}
\newcommand{\AnnotationTok}[1]{\textcolor[rgb]{0.56,0.35,0.01}{\textbf{\textit{#1}}}}
\newcommand{\AttributeTok}[1]{\textcolor[rgb]{0.13,0.29,0.53}{#1}}
\newcommand{\BaseNTok}[1]{\textcolor[rgb]{0.00,0.00,0.81}{#1}}
\newcommand{\BuiltInTok}[1]{#1}
\newcommand{\CharTok}[1]{\textcolor[rgb]{0.31,0.60,0.02}{#1}}
\newcommand{\CommentTok}[1]{\textcolor[rgb]{0.56,0.35,0.01}{\textit{#1}}}
\newcommand{\CommentVarTok}[1]{\textcolor[rgb]{0.56,0.35,0.01}{\textbf{\textit{#1}}}}
\newcommand{\ConstantTok}[1]{\textcolor[rgb]{0.56,0.35,0.01}{#1}}
\newcommand{\ControlFlowTok}[1]{\textcolor[rgb]{0.13,0.29,0.53}{\textbf{#1}}}
\newcommand{\DataTypeTok}[1]{\textcolor[rgb]{0.13,0.29,0.53}{#1}}
\newcommand{\DecValTok}[1]{\textcolor[rgb]{0.00,0.00,0.81}{#1}}
\newcommand{\DocumentationTok}[1]{\textcolor[rgb]{0.56,0.35,0.01}{\textbf{\textit{#1}}}}
\newcommand{\ErrorTok}[1]{\textcolor[rgb]{0.64,0.00,0.00}{\textbf{#1}}}
\newcommand{\ExtensionTok}[1]{#1}
\newcommand{\FloatTok}[1]{\textcolor[rgb]{0.00,0.00,0.81}{#1}}
\newcommand{\FunctionTok}[1]{\textcolor[rgb]{0.13,0.29,0.53}{\textbf{#1}}}
\newcommand{\ImportTok}[1]{#1}
\newcommand{\InformationTok}[1]{\textcolor[rgb]{0.56,0.35,0.01}{\textbf{\textit{#1}}}}
\newcommand{\KeywordTok}[1]{\textcolor[rgb]{0.13,0.29,0.53}{\textbf{#1}}}
\newcommand{\NormalTok}[1]{#1}
\newcommand{\OperatorTok}[1]{\textcolor[rgb]{0.81,0.36,0.00}{\textbf{#1}}}
\newcommand{\OtherTok}[1]{\textcolor[rgb]{0.56,0.35,0.01}{#1}}
\newcommand{\PreprocessorTok}[1]{\textcolor[rgb]{0.56,0.35,0.01}{\textit{#1}}}
\newcommand{\RegionMarkerTok}[1]{#1}
\newcommand{\SpecialCharTok}[1]{\textcolor[rgb]{0.81,0.36,0.00}{\textbf{#1}}}
\newcommand{\SpecialStringTok}[1]{\textcolor[rgb]{0.31,0.60,0.02}{#1}}
\newcommand{\StringTok}[1]{\textcolor[rgb]{0.31,0.60,0.02}{#1}}
\newcommand{\VariableTok}[1]{\textcolor[rgb]{0.00,0.00,0.00}{#1}}
\newcommand{\VerbatimStringTok}[1]{\textcolor[rgb]{0.31,0.60,0.02}{#1}}
\newcommand{\WarningTok}[1]{\textcolor[rgb]{0.56,0.35,0.01}{\textbf{\textit{#1}}}}
\usepackage{longtable,booktabs,array}
\usepackage{calc} % for calculating minipage widths
% Correct order of tables after \paragraph or \subparagraph
\usepackage{etoolbox}
\makeatletter
\patchcmd\longtable{\par}{\if@noskipsec\mbox{}\fi\par}{}{}
\makeatother
% Allow footnotes in longtable head/foot
\IfFileExists{footnotehyper.sty}{\usepackage{footnotehyper}}{\usepackage{footnote}}
\makesavenoteenv{longtable}
\usepackage{graphicx}
\makeatletter
\newsavebox\pandoc@box
\newcommand*\pandocbounded[1]{% scales image to fit in text height/width
  \sbox\pandoc@box{#1}%
  \Gscale@div\@tempa{\textheight}{\dimexpr\ht\pandoc@box+\dp\pandoc@box\relax}%
  \Gscale@div\@tempb{\linewidth}{\wd\pandoc@box}%
  \ifdim\@tempb\p@<\@tempa\p@\let\@tempa\@tempb\fi% select the smaller of both
  \ifdim\@tempa\p@<\p@\scalebox{\@tempa}{\usebox\pandoc@box}%
  \else\usebox{\pandoc@box}%
  \fi%
}
% Set default figure placement to htbp
\def\fps@figure{htbp}
\makeatother
\setlength{\emergencystretch}{3em} % prevent overfull lines
\providecommand{\tightlist}{%
  \setlength{\itemsep}{0pt}\setlength{\parskip}{0pt}}
\usepackage{bookmark}
\IfFileExists{xurl.sty}{\usepackage{xurl}}{} % add URL line breaks if available
\urlstyle{same}
\hypersetup{
  pdftitle={CHULETA DECISIÓN - Data Science y Diseño Experimental},
  pdfauthor={Guía Rápida para Examen},
  hidelinks,
  pdfcreator={LaTeX via pandoc}}

\title{CHULETA DECISIÓN - Data Science y Diseño Experimental}
\usepackage{etoolbox}
\makeatletter
\providecommand{\subtitle}[1]{% add subtitle to \maketitle
  \apptocmd{\@title}{\par {\large #1 \par}}{}{}
}
\makeatother
\subtitle{¿Qué Test/Método Usar y Cuándo?}
\author{Guía Rápida para Examen}
\date{2025-10-14}

\begin{document}
\maketitle

{
\setcounter{tocdepth}{3}
\tableofcontents
}
\newpage

\section{ANÁLISIS EXPLORATORIO
INICIAL}\label{anuxe1lisis-exploratorio-inicial}

\subsection{¿Qué tipo de datos tengo?}\label{quuxe9-tipo-de-datos-tengo}

\subsubsection{VARIABLES CONTINUAS}\label{variables-continuas}

Números con decimales: peso, altura, precio

\textbf{Ejemplo pregunta}: \emph{``Analiza la distribución del peso de
los pacientes''}

\begin{itemize}
\tightlist
\item
  → Descriptivos: media, mediana, sd, histogramas, boxplots
\item
  → Verificar: normalidad (Shapiro-Wilk), outliers
\end{itemize}

\subsubsection{VARIABLES CATEGÓRICAS}\label{variables-categuxf3ricas}

Grupos: género, ciudad, sí/no

\textbf{Ejemplo pregunta}: \emph{``¿Cuántos pacientes hay por género?''}

\begin{itemize}
\tightlist
\item
  → Descriptivos: tablas de frecuencia, barplots
\item
  → Verificar: distribución de categorías
\end{itemize}

\subsubsection{VARIABLES DE CONTEO}\label{variables-de-conteo}

Enteros ≥ 0: número de visitas, eventos

\textbf{Ejemplo pregunta}: \emph{``Describe el número de visitas
hospitalarias''}

\begin{itemize}
\tightlist
\item
  → Descriptivos: media, rango, histogramas
\end{itemize}

\newpage

\section{TESTS DE HIPÓTESIS - ÁRBOL DE
DECISIÓN}\label{tests-de-hipuxf3tesis---uxe1rbol-de-decisiuxf3n}

\subsection{VARIABLES CONTINUAS
(medias/localizaciones)}\label{variables-continuas-mediaslocalizaciones}

\subsubsection{1 GRUPO (comparar con valor de
referencia)}\label{grupo-comparar-con-valor-de-referencia}

\textbf{Ejemplo pregunta}: \emph{``¿La altura media de los estudiantes
es diferente de 170cm?''}

\begin{itemize}
\tightlist
\item
  ✅ \textbf{One-sample t-test} si: datos normales, sin outliers
\item
  ✅ \textbf{Wilcoxon signed-rank} si: outliers o no normalidad
\end{itemize}

\begin{Shaded}
\begin{Highlighting}[]
\CommentTok{\# Si normal:}
\FunctionTok{t.test}\NormalTok{(alturas, }\AttributeTok{mu =} \DecValTok{170}\NormalTok{)}

\CommentTok{\# Si no normal:}
\FunctionTok{wilcox.test}\NormalTok{(alturas, }\AttributeTok{mu =} \DecValTok{170}\NormalTok{)}
\end{Highlighting}
\end{Shaded}

\subsubsection{2 GRUPOS - RELACIONADAS}\label{grupos---relacionadas}

\textbf{Ejemplo pregunta}: \emph{``¿El peso de los pacientes cambió
después del tratamiento?''}

\textbf{Ejemplo pregunta}: \emph{``¿Hay diferencia en la presión
arterial antes y después del ejercicio?''}

\begin{itemize}
\tightlist
\item
  ✅ \textbf{Paired t-test} si: diferencias normales
\item
  ✅ \textbf{Wilcoxon signed-rank} si: outliers o no normalidad
\end{itemize}

\begin{Shaded}
\begin{Highlighting}[]
\CommentTok{\# Si normal:}
\FunctionTok{t.test}\NormalTok{(peso\_antes, peso\_despues, }\AttributeTok{paired =} \ConstantTok{TRUE}\NormalTok{)}

\CommentTok{\# Si no normal:}
\FunctionTok{wilcox.test}\NormalTok{(peso\_antes, peso\_despues, }\AttributeTok{paired =} \ConstantTok{TRUE}\NormalTok{)}
\end{Highlighting}
\end{Shaded}

\subsubsection{2 GRUPOS - INDEPENDIENTES}\label{grupos---independientes}

\textbf{Ejemplo pregunta}: \emph{``¿Los hombres pesan más que las
mujeres?''}

\textbf{Ejemplo pregunta}: \emph{``¿Hay diferencia en el salario entre
fumadores y no fumadores?''}

\textbf{¿Datos normales y sin outliers?}

\begin{itemize}
\tightlist
\item
  \textbf{SÍ} → Two-sample t-test
\item
  \textbf{NO} → Wilcoxon-Mann-Whitney
\end{itemize}

\begin{Shaded}
\begin{Highlighting}[]
\CommentTok{\# Verificar normalidad primero:}
\FunctionTok{shapiro.test}\NormalTok{(grupo1)}
\FunctionTok{shapiro.test}\NormalTok{(grupo2)}

\CommentTok{\# Verificar varianzas:}
\FunctionTok{var.test}\NormalTok{(grupo1, grupo2)}

\CommentTok{\# Si varianzas iguales:}
\FunctionTok{t.test}\NormalTok{(grupo1, grupo2, }\AttributeTok{var.equal =} \ConstantTok{TRUE}\NormalTok{)}

\CommentTok{\# Si varianzas diferentes:}
\FunctionTok{t.test}\NormalTok{(grupo1, grupo2, }\AttributeTok{var.equal =} \ConstantTok{FALSE}\NormalTok{)}

\CommentTok{\# Si no normal:}
\FunctionTok{wilcox.test}\NormalTok{(grupo1, grupo2)}
\end{Highlighting}
\end{Shaded}

\subsubsection{3+ GRUPOS independientes}\label{grupos-independientes}

\textbf{Ejemplo pregunta}: \emph{``¿El peso varía entre personas de
Asia, Europa y América?''}

\textbf{Ejemplo pregunta}: \emph{``Compara el rendimiento académico
entre 4 escuelas diferentes''}

\textbf{¿Datos normales, varianzas homogéneas, sin outliers?}

\begin{itemize}
\tightlist
\item
  \textbf{SÍ} → ✅ ANOVA
\item
  \textbf{NO} → ✅ Kruskal-Wallis
\end{itemize}

\begin{Shaded}
\begin{Highlighting}[]
\CommentTok{\# Verificar:}
\FunctionTok{shapiro.test}\NormalTok{(datos) }\CommentTok{\# para cada grupo}
\FunctionTok{bartlett.test}\NormalTok{(variable }\SpecialCharTok{\textasciitilde{}}\NormalTok{ grupo)}

\CommentTok{\# ANOVA:}
\NormalTok{modelo }\OtherTok{\textless{}{-}} \FunctionTok{aov}\NormalTok{(variable }\SpecialCharTok{\textasciitilde{}}\NormalTok{ grupo)}
\FunctionTok{summary}\NormalTok{(modelo)}

\CommentTok{\# Si no paramétrico:}
\FunctionTok{kruskal.test}\NormalTok{(variable }\SpecialCharTok{\textasciitilde{}}\NormalTok{ grupo)}
\end{Highlighting}
\end{Shaded}

⚠️ \textbf{IMPORTANTE}: ANOVA/Kruskal-Wallis dicen ``hay diferencias''
pero NO cuál grupo

\newpage

\subsection{VARIANZAS/DISPERSIONES}\label{varianzasdispersiones}

\subsubsection{2 GRUPOS}\label{grupos}

\textbf{Ejemplo pregunta}: \emph{``¿La variabilidad del peso es la misma
en hombres y mujeres?''}

\begin{Shaded}
\begin{Highlighting}[]
\FunctionTok{var.test}\NormalTok{(grupo1, grupo2)}
\end{Highlighting}
\end{Shaded}

\subsubsection{3+ GRUPOS}\label{grupos-1}

\textbf{Ejemplo pregunta}: \emph{``¿La dispersión de las notas es
homogénea entre las 5 clases?''}

\begin{Shaded}
\begin{Highlighting}[]
\FunctionTok{bartlett.test}\NormalTok{(notas }\SpecialCharTok{\textasciitilde{}}\NormalTok{ clase)}
\end{Highlighting}
\end{Shaded}

\subsection{VARIABLES CATEGÓRICAS
(asociación/independencia)}\label{variables-categuxf3ricas-asociaciuxf3nindependencia}

\subsubsection{TABLA 2×2}\label{tabla-22}

\textbf{Ejemplo pregunta}: \emph{``¿Ser catalán está asociado con ser
del Barça?''}

\textbf{Ejemplo pregunta}: \emph{``¿Hay asociación entre fumar y tener
cáncer de pulmón?''}

\textbf{Ejemplo pregunta}: \emph{``¿Jugar en casa afecta a las victorias
del Real Sociedad?''}

\textbf{¿Muestra pequeña (frecuencias esperadas \textless{} 5)?}

\begin{itemize}
\tightlist
\item
  \textbf{SÍ} → ✅ Test de Fisher (exacto)
\item
  \textbf{NO} → ✅ Test de Chi-cuadrado
\end{itemize}

\begin{Shaded}
\begin{Highlighting}[]
\CommentTok{\# Fisher (muestras pequeñas):}
\FunctionTok{fisher.test}\NormalTok{(tabla, }\AttributeTok{alternative =} \StringTok{"greater"}\NormalTok{)}

\CommentTok{\# Chi{-}cuadrado (muestras grandes):}
\FunctionTok{chisq.test}\NormalTok{(tabla)}
\end{Highlighting}
\end{Shaded}

\subsubsection{TABLA mayor que 2×2}\label{tabla-mayor-que-22}

\textbf{Ejemplo pregunta}: \emph{``¿Hay asociación entre estado y
población?''}

\begin{Shaded}
\begin{Highlighting}[]
\FunctionTok{chisq.test}\NormalTok{(tabla)}
\end{Highlighting}
\end{Shaded}

\subsubsection{PROPORCIÓN BINARIA}\label{proporciuxf3n-binaria}

\textbf{Ejemplo pregunta}: \emph{``Paul el pulpo acertó 8 de 8 partidos.
¿Es suerte?''}

\textbf{Ejemplo pregunta}: \emph{``Una moneda salió cara 8 de 10 veces.
¿Está trucada?''}

\begin{Shaded}
\begin{Highlighting}[]
\FunctionTok{binom.test}\NormalTok{(}\DecValTok{8}\NormalTok{, }\DecValTok{10}\NormalTok{, }\AttributeTok{p =} \FloatTok{0.5}\NormalTok{, }\AttributeTok{alternative =} \StringTok{"greater"}\NormalTok{)}
\FunctionTok{binom.test}\NormalTok{(}\DecValTok{8}\NormalTok{, }\DecValTok{8}\NormalTok{, }\AttributeTok{p =} \FloatTok{0.5}\NormalTok{, }\AttributeTok{alternative =} \StringTok{"greater"}\NormalTok{)}
\end{Highlighting}
\end{Shaded}

\subsection{CUANDO NO SÉ LA DISTRIBUCIÓN /
DUDAS}\label{cuando-no-suxe9-la-distribuciuxf3n-dudas}

\subsubsection{Permutation Test}\label{permutation-test}

\textbf{Ejemplo pregunta}: \emph{``Compara dos grupos con outliers y
dudas sobre test''}

\begin{Shaded}
\begin{Highlighting}[]
\CommentTok{\# Creas tu propia distribución nula permutando etiquetas}
\NormalTok{diferencia\_observada }\OtherTok{\textless{}{-}} \FunctionTok{mean}\NormalTok{(grupo\_B) }\SpecialCharTok{{-}} \FunctionTok{mean}\NormalTok{(grupo\_A)}
\NormalTok{datos\_combinados }\OtherTok{\textless{}{-}} \FunctionTok{c}\NormalTok{(grupo\_A, grupo\_B)}

\NormalTok{diferencias\_simuladas }\OtherTok{\textless{}{-}} \FunctionTok{numeric}\NormalTok{(}\DecValTok{10000}\NormalTok{)}
\ControlFlowTok{for}\NormalTok{ (i }\ControlFlowTok{in} \DecValTok{1}\SpecialCharTok{:}\DecValTok{10000}\NormalTok{) \{}
\NormalTok{  datos\_permutados }\OtherTok{\textless{}{-}} \FunctionTok{sample}\NormalTok{(datos\_combinados)}
\NormalTok{  grupo\_A\_sim }\OtherTok{\textless{}{-}}\NormalTok{ datos\_permutados[}\DecValTok{1}\SpecialCharTok{:}\FunctionTok{length}\NormalTok{(grupo\_A)]}
\NormalTok{  grupo\_B\_sim }\OtherTok{\textless{}{-}}\NormalTok{ datos\_permutados[(}\FunctionTok{length}\NormalTok{(grupo\_A)}\SpecialCharTok{+}\DecValTok{1}\NormalTok{)}\SpecialCharTok{:}\FunctionTok{length}\NormalTok{(datos\_combinados)]}
\NormalTok{  diferencias\_simuladas[i] }\OtherTok{\textless{}{-}} \FunctionTok{mean}\NormalTok{(grupo\_B\_sim) }\SpecialCharTok{{-}} \FunctionTok{mean}\NormalTok{(grupo\_A\_sim)}
\NormalTok{\}}

\NormalTok{p\_value }\OtherTok{\textless{}{-}} \FunctionTok{sum}\NormalTok{(}\FunctionTok{abs}\NormalTok{(diferencias\_simuladas) }\SpecialCharTok{\textgreater{}=} \FunctionTok{abs}\NormalTok{(diferencia\_observada)) }\SpecialCharTok{/} \DecValTok{10000}
\end{Highlighting}
\end{Shaded}

\subsubsection{Bootstrap}\label{bootstrap}

\textbf{Ejemplo pregunta}: \emph{``¿Cuál es el intervalo de confianza de
la mediana?''}

\begin{Shaded}
\begin{Highlighting}[]
\CommentTok{\# Remuestreo con reemplazo para estimar distribución}
\NormalTok{medianas\_bootstrap }\OtherTok{\textless{}{-}} \FunctionTok{numeric}\NormalTok{(}\DecValTok{10000}\NormalTok{)}
\ControlFlowTok{for}\NormalTok{(i }\ControlFlowTok{in} \DecValTok{1}\SpecialCharTok{:}\DecValTok{10000}\NormalTok{)\{}
\NormalTok{  medianas\_bootstrap[i] }\OtherTok{\textless{}{-}} \FunctionTok{median}\NormalTok{(}\FunctionTok{sample}\NormalTok{(x, }\AttributeTok{size=}\FunctionTok{length}\NormalTok{(x), }\AttributeTok{replace=}\ConstantTok{TRUE}\NormalTok{))}
\NormalTok{\}}
\FunctionTok{quantile}\NormalTok{(medianas\_bootstrap, }\FunctionTok{c}\NormalTok{(}\FloatTok{0.025}\NormalTok{, }\FloatTok{0.975}\NormalTok{))}
\end{Highlighting}
\end{Shaded}

\newpage

\section{VERIFICACIÓN DE SUPUESTOS}\label{verificaciuxf3n-de-supuestos}

\subsection{¿Mis datos son normales?}\label{mis-datos-son-normales}

\textbf{Ejemplo pregunta}: \emph{``Antes de hacer un t-test, verifica
normalidad''}

\begin{Shaded}
\begin{Highlighting}[]
\FunctionTok{shapiro.test}\NormalTok{(datos)}
\CommentTok{\# H₀: datos son normales}
\CommentTok{\# Si p \textless{} 0.05 → NO normales → USA }\AlertTok{TEST}\CommentTok{ NO PARAMÉTRICO}
\end{Highlighting}
\end{Shaded}

\subsection{¿Hay outliers extremos?}\label{hay-outliers-extremos}

\textbf{Ejemplo pregunta}: \emph{``Identifica valores atípicos en el
dataset''}

\begin{Shaded}
\begin{Highlighting}[]
\FunctionTok{boxplot}\NormalTok{(datos)}
\CommentTok{\# Si hay muchos outliers → USA }\AlertTok{TEST}\CommentTok{ NO PARAMÉTRICO}
\end{Highlighting}
\end{Shaded}

\subsection{¿Varianzas homogéneas?}\label{varianzas-homoguxe9neas}

\textbf{Ejemplo pregunta}: \emph{``Verifica si puedes usar
var.equal=TRUE''}

\begin{Shaded}
\begin{Highlighting}[]
\CommentTok{\# 2 grupos:}
\FunctionTok{var.test}\NormalTok{(grupo1, grupo2)}

\CommentTok{\# 3+ grupos:}
\FunctionTok{bartlett.test}\NormalTok{(variable }\SpecialCharTok{\textasciitilde{}}\NormalTok{ grupo)}

\CommentTok{\# Si p \textless{} 0.05 → varianzas diferentes}
\CommentTok{\# En t{-}test: usa var.equal = FALSE}
\end{Highlighting}
\end{Shaded}

\newpage

\section{REGRESIÓN - ÁRBOL DE
DECISIÓN}\label{regresiuxf3n---uxe1rbol-de-decisiuxf3n}

\subsection{Y es CONTINUA (peso, precio,
temperatura)}\label{y-es-continua-peso-precio-temperatura}

\subsubsection{1 PREDICTOR}\label{predictor}

\textbf{Ejemplo pregunta}: \emph{``¿Cómo afecta el peso del coche al
consumo (mpg)?''}

\textbf{Ejemplo pregunta}: \emph{``¿El BMI predice los gastos
médicos?''}

\begin{Shaded}
\begin{Highlighting}[]
\NormalTok{modelo }\OtherTok{\textless{}{-}} \FunctionTok{lm}\NormalTok{(mpg }\SpecialCharTok{\textasciitilde{}}\NormalTok{ weight, }\AttributeTok{data =}\NormalTok{ datos)}
\FunctionTok{summary}\NormalTok{(modelo)}
\end{Highlighting}
\end{Shaded}

\subsubsection{MÚLTIPLES PREDICTORES}\label{muxfaltiples-predictores}

\textbf{Ejemplo pregunta}: \emph{``¿Cómo afectan edad, BMI y fumar a los
gastos médicos?''}

\textbf{Ejemplo pregunta}: \emph{``Predice el precio de venta de una
casa usando múltiples características''}

\begin{Shaded}
\begin{Highlighting}[]
\NormalTok{modelo }\OtherTok{\textless{}{-}} \FunctionTok{lm}\NormalTok{(charges }\SpecialCharTok{\textasciitilde{}}\NormalTok{ age }\SpecialCharTok{+}\NormalTok{ bmi }\SpecialCharTok{+}\NormalTok{ smoker, }\AttributeTok{data =}\NormalTok{ insurance)}
\FunctionTok{summary}\NormalTok{(modelo)}
\end{Highlighting}
\end{Shaded}

\subsubsection{VERIFICAR SUPUESTOS}\label{verificar-supuestos}

\begin{Shaded}
\begin{Highlighting}[]
\CommentTok{\# 1. Linealidad}
\FunctionTok{plot}\NormalTok{(datos}\SpecialCharTok{$}\NormalTok{Y }\SpecialCharTok{\textasciitilde{}}\NormalTok{ datos}\SpecialCharTok{$}\NormalTok{X)}

\CommentTok{\# 2. Normalidad residuos}
\FunctionTok{shapiro.test}\NormalTok{(}\FunctionTok{residuals}\NormalTok{(modelo))}

\CommentTok{\# 3. Homocedasticidad}
\FunctionTok{plot}\NormalTok{(modelo, }\AttributeTok{which =} \DecValTok{1}\NormalTok{)}

\CommentTok{\# 4. Outliers influyentes}
\FunctionTok{plot}\NormalTok{(modelo, }\AttributeTok{which =} \DecValTok{5}\NormalTok{)}
\end{Highlighting}
\end{Shaded}

\subsubsection{SI NO SE CUMPLEN
SUPUESTOS}\label{si-no-se-cumplen-supuestos}

\begin{Shaded}
\begin{Highlighting}[]
\CommentTok{\# Transformación:}
\NormalTok{modelo }\OtherTok{\textless{}{-}} \FunctionTok{lm}\NormalTok{(}\FunctionTok{log}\NormalTok{(Y) }\SpecialCharTok{\textasciitilde{}}\NormalTok{ X, }\AttributeTok{data =}\NormalTok{ datos)}
\NormalTok{modelo }\OtherTok{\textless{}{-}} \FunctionTok{lm}\NormalTok{(}\FunctionTok{sqrt}\NormalTok{(Y) }\SpecialCharTok{\textasciitilde{}}\NormalTok{ X, }\AttributeTok{data =}\NormalTok{ datos)}
\end{Highlighting}
\end{Shaded}

\subsection{Y es BINARIA (0/1, sí/no,
enfermo/sano)}\label{y-es-binaria-01-suxedno-enfermosano}

\textbf{Ejemplo pregunta}: \emph{``¿El BMI predice si un paciente tendrá
un evento cardíaco?''}

\textbf{Ejemplo pregunta}: \emph{``¿La edad predice si alguien aprueba
el examen?''}

\begin{Shaded}
\begin{Highlighting}[]
\NormalTok{modelo }\OtherTok{\textless{}{-}} \FunctionTok{glm}\NormalTok{(evento\_cardiaco }\SpecialCharTok{\textasciitilde{}}\NormalTok{ BMI, }
              \AttributeTok{family =} \FunctionTok{binomial}\NormalTok{(}\AttributeTok{link =} \StringTok{"logit"}\NormalTok{),}
              \AttributeTok{data =}\NormalTok{ datos)}
\FunctionTok{summary}\NormalTok{(modelo)}
\end{Highlighting}
\end{Shaded}

\textbf{Distribución}: Binomial\\
\textbf{Link}: logit\\
\textbf{Y debe estar entre 0 y 1} (probabilidad)

\subsection{Y es CONTEO (número entero ≥
0)}\label{y-es-conteo-nuxfamero-entero-0}

\textbf{Ejemplo pregunta}: \emph{``¿El BMI predice el número de visitas
hospitalarias?''}

\textbf{Ejemplo pregunta}: \emph{``Modela el número de accidentes según
variables de tráfico''}

\begin{Shaded}
\begin{Highlighting}[]
\NormalTok{modelo }\OtherTok{\textless{}{-}} \FunctionTok{glm}\NormalTok{(num\_visitas }\SpecialCharTok{\textasciitilde{}}\NormalTok{ BMI, }
              \AttributeTok{family =} \FunctionTok{poisson}\NormalTok{(}\AttributeTok{link =} \StringTok{"log"}\NormalTok{),}
              \AttributeTok{data =}\NormalTok{ datos)}
\FunctionTok{summary}\NormalTok{(modelo)}
\end{Highlighting}
\end{Shaded}

\textbf{Distribución}: Poisson\\
\textbf{Link}: log\\
\textbf{Y debe ser ≥ 0} (conteos)

\subsection{Y CONTINUA con PROBLEMAS}\label{y-continua-con-problemas}

\textbf{Ejemplo pregunta}: \emph{``BMI vs gastos médicos tiene relación
no lineal''}

\subsubsection{Opción 1: Transformar Y}\label{opciuxf3n-1-transformar-y}

\begin{Shaded}
\begin{Highlighting}[]
\NormalTok{modelo }\OtherTok{\textless{}{-}} \FunctionTok{lm}\NormalTok{(}\FunctionTok{log}\NormalTok{(charges) }\SpecialCharTok{\textasciitilde{}}\NormalTok{ BMI, }\AttributeTok{data =}\NormalTok{ datos)}
\NormalTok{modelo }\OtherTok{\textless{}{-}} \FunctionTok{lm}\NormalTok{(}\FunctionTok{sqrt}\NormalTok{(charges) }\SpecialCharTok{\textasciitilde{}}\NormalTok{ BMI, }\AttributeTok{data =}\NormalTok{ datos)}
\end{Highlighting}
\end{Shaded}

\subsubsection{Opción 2: GLM con Gamma}\label{opciuxf3n-2-glm-con-gamma}

\begin{Shaded}
\begin{Highlighting}[]
\NormalTok{modelo }\OtherTok{\textless{}{-}} \FunctionTok{glm}\NormalTok{(charges }\SpecialCharTok{\textasciitilde{}}\NormalTok{ BMI, }
              \AttributeTok{family =} \FunctionTok{Gamma}\NormalTok{(}\AttributeTok{link =} \StringTok{"log"}\NormalTok{),}
              \AttributeTok{data =}\NormalTok{ datos)}
\end{Highlighting}
\end{Shaded}

\newpage

\section{SCREENING DE VARIABLES}\label{screening-de-variables}

\subsection{¿Tengo 10+ variables
predictoras?}\label{tengo-10-variables-predictoras}

\textbf{Ejemplo pregunta}: \emph{``Dataset con 80 predictores,
identifica los significativos''}

\textbf{PROBLEMA}: Muchos tests → \textasciitilde5\% falsos positivos
esperados

\subsubsection{Corrección FDR
(Benjamini-Hochberg)}\label{correcciuxf3n-fdr-benjamini-hochberg}

✅ \textbf{USA ESTO en análisis exploratorios}

\begin{Shaded}
\begin{Highlighting}[]
\CommentTok{\# Hacer muchas regresiones y guardar p{-}values}
\NormalTok{p\_values }\OtherTok{\textless{}{-}} \FunctionTok{numeric}\NormalTok{(}\FunctionTok{length}\NormalTok{(predictores))}
\ControlFlowTok{for}\NormalTok{(i }\ControlFlowTok{in} \DecValTok{1}\SpecialCharTok{:}\FunctionTok{length}\NormalTok{(predictores)) \{}
\NormalTok{  modelo }\OtherTok{\textless{}{-}} \FunctionTok{lm}\NormalTok{(Y }\SpecialCharTok{\textasciitilde{}}\NormalTok{ datos[,predictores[i]])}
\NormalTok{  p\_values[i] }\OtherTok{\textless{}{-}} \FunctionTok{summary}\NormalTok{(modelo)}\SpecialCharTok{$}\NormalTok{coefficients[}\DecValTok{2}\NormalTok{,}\DecValTok{4}\NormalTok{]}
\NormalTok{\}}

\CommentTok{\# Corrección FDR}
\NormalTok{p\_adjusted }\OtherTok{\textless{}{-}} \FunctionTok{p.adjust}\NormalTok{(p\_values, }\AttributeTok{method =} \StringTok{"BH"}\NormalTok{)}

\CommentTok{\# Ver cuáles siguen siendo significativos}
\NormalTok{significant }\OtherTok{\textless{}{-}} \FunctionTok{which}\NormalTok{(p\_adjusted }\SpecialCharTok{\textless{}} \FloatTok{0.05}\NormalTok{)}
\end{Highlighting}
\end{Shaded}

\subsubsection{Corrección Bonferroni}\label{correcciuxf3n-bonferroni}

Solo si DEBES evitar TODO falso positivo

\begin{Shaded}
\begin{Highlighting}[]
\NormalTok{p\_adjusted }\OtherTok{\textless{}{-}} \FunctionTok{p.adjust}\NormalTok{(p\_values, }\AttributeTok{method =} \StringTok{"bonferroni"}\NormalTok{)}
\end{Highlighting}
\end{Shaded}

\newpage

\section{REDUCCIÓN DE
DIMENSIONALIDAD}\label{reducciuxf3n-de-dimensionalidad}

\subsection{VISUALIZACIÓN}\label{visualizaciuxf3n}

\textbf{Ejemplo pregunta}: \emph{``Tienes datos de 60 sensores,
visualiza patrones''}

\textbf{Ejemplo pregunta}: \emph{``Visualiza jugadores de fútbol según
30+ atributos''}

\begin{Shaded}
\begin{Highlighting}[]
\CommentTok{\# PCA}
\NormalTok{pca\_result }\OtherTok{\textless{}{-}} \FunctionTok{prcomp}\NormalTok{(datos, }\AttributeTok{center =} \ConstantTok{TRUE}\NormalTok{, }\AttributeTok{scale. =} \ConstantTok{TRUE}\NormalTok{)}

\CommentTok{\# Ver varianza explicada}
\FunctionTok{summary}\NormalTok{(pca\_result)}

\CommentTok{\# Visualizar}
\FunctionTok{plot}\NormalTok{(pca\_result}\SpecialCharTok{$}\NormalTok{x[,}\DecValTok{1}\NormalTok{], pca\_result}\SpecialCharTok{$}\NormalTok{x[,}\DecValTok{2}\NormalTok{], }
     \AttributeTok{col =}\NormalTok{ grupos, }\AttributeTok{pch =} \DecValTok{19}\NormalTok{,}
     \AttributeTok{xlab =} \StringTok{"PC1"}\NormalTok{, }\AttributeTok{ylab =} \StringTok{"PC2"}\NormalTok{)}
\FunctionTok{legend}\NormalTok{(}\StringTok{"topright"}\NormalTok{, }\AttributeTok{legend =} \FunctionTok{levels}\NormalTok{(grupos), }
       \AttributeTok{col =} \DecValTok{1}\SpecialCharTok{:}\FunctionTok{length}\NormalTok{(}\FunctionTok{levels}\NormalTok{(grupos)), }\AttributeTok{pch =} \DecValTok{19}\NormalTok{)}

\CommentTok{\# Con ggplot}
\FunctionTok{library}\NormalTok{(ggplot2)}
\NormalTok{data\_pca }\OtherTok{\textless{}{-}} \FunctionTok{data.frame}\NormalTok{(}\AttributeTok{PC1 =}\NormalTok{ pca\_result}\SpecialCharTok{$}\NormalTok{x[,}\DecValTok{1}\NormalTok{],}
                       \AttributeTok{PC2 =}\NormalTok{ pca\_result}\SpecialCharTok{$}\NormalTok{x[,}\DecValTok{2}\NormalTok{],}
                       \AttributeTok{grupo =}\NormalTok{ grupos)}
\FunctionTok{ggplot}\NormalTok{(data\_pca, }\FunctionTok{aes}\NormalTok{(PC1, PC2, }\AttributeTok{color =}\NormalTok{ grupo)) }\SpecialCharTok{+}
  \FunctionTok{geom\_point}\NormalTok{() }\SpecialCharTok{+}
  \FunctionTok{theme\_minimal}\NormalTok{()}
\end{Highlighting}
\end{Shaded}

\subsubsection{¿Cuántos componentes
usar?}\label{cuuxe1ntos-componentes-usar}

\begin{Shaded}
\begin{Highlighting}[]
\CommentTok{\# Scree plot}
\FunctionTok{plot}\NormalTok{(pca\_result, }\AttributeTok{type =} \StringTok{"l"}\NormalTok{)}

\CommentTok{\# Varianza explicada acumulada}
\FunctionTok{cumsum}\NormalTok{(pca\_result}\SpecialCharTok{$}\NormalTok{sdev}\SpecialCharTok{\^{}}\DecValTok{2} \SpecialCharTok{/} \FunctionTok{sum}\NormalTok{(pca\_result}\SpecialCharTok{$}\NormalTok{sdev}\SpecialCharTok{\^{}}\DecValTok{2}\NormalTok{))}
\CommentTok{\# Típicamente: 80{-}90\% de varianza}
\end{Highlighting}
\end{Shaded}

\subsection{ELIMINAR
MULTICOLINEALIDAD}\label{eliminar-multicolinealidad}

\textbf{Ejemplo pregunta}: \emph{``20 predictores correlacionados antes
de regresión''}

\begin{Shaded}
\begin{Highlighting}[]
\NormalTok{pca\_result }\OtherTok{\textless{}{-}} \FunctionTok{prcomp}\NormalTok{(predictores, }\AttributeTok{center =} \ConstantTok{TRUE}\NormalTok{, }\AttributeTok{scale. =} \ConstantTok{TRUE}\NormalTok{)}
\CommentTok{\# Usa los primeros n componentes}
\NormalTok{datos\_pca }\OtherTok{\textless{}{-}}\NormalTok{ pca\_result}\SpecialCharTok{$}\NormalTok{x[, }\DecValTok{1}\SpecialCharTok{:}\DecValTok{5}\NormalTok{]}
\NormalTok{modelo }\OtherTok{\textless{}{-}} \FunctionTok{lm}\NormalTok{(Y }\SpecialCharTok{\textasciitilde{}}\NormalTok{ datos\_pca)}
\end{Highlighting}
\end{Shaded}

\newpage

\section{DISEÑO EXPERIMENTAL}\label{diseuxf1o-experimental}

\subsection{¿Cuántas muestras
necesito?}\label{cuuxe1ntas-muestras-necesito}

\textbf{Ejemplo pregunta}: \emph{``¿Cuántos pacientes necesitas para
comparar dos tratamientos?''}

\begin{Shaded}
\begin{Highlighting}[]
\FunctionTok{library}\NormalTok{(pwr)}

\CommentTok{\# Para t{-}test}
\FunctionTok{pwr.t.test}\NormalTok{(}\AttributeTok{d =} \FloatTok{0.5}\NormalTok{,           }\CommentTok{\# effect size}
           \AttributeTok{sig.level =} \FloatTok{0.05}\NormalTok{,   }\CommentTok{\# alfa}
           \AttributeTok{power =} \FloatTok{0.8}\NormalTok{,        }\CommentTok{\# poder deseado}
           \AttributeTok{type =} \StringTok{"two.sample"}\NormalTok{)}

\CommentTok{\# Para ANOVA}
\FunctionTok{pwr.anova.test}\NormalTok{(}\AttributeTok{k =} \DecValTok{3}\NormalTok{,          }\CommentTok{\# número de grupos}
               \AttributeTok{f =} \FloatTok{0.25}\NormalTok{,       }\CommentTok{\# effect size}
               \AttributeTok{sig.level =} \FloatTok{0.05}\NormalTok{,}
               \AttributeTok{power =} \FloatTok{0.8}\NormalTok{)}
\end{Highlighting}
\end{Shaded}

\subsection{¿Tengo una fuente conocida de
variabilidad?}\label{tengo-una-fuente-conocida-de-variabilidad}

\textbf{Ejemplo pregunta}: \emph{``Comparas 3 dietas pero el género
afecta al peso''}

\textbf{Ejemplo pregunta}: \emph{``Experimento con zapatos: cada persona
prueba ambos materiales''}

\begin{Shaded}
\begin{Highlighting}[]
\CommentTok{\# Diseño de bloques}
\NormalTok{modelo }\OtherTok{\textless{}{-}} \FunctionTok{aov}\NormalTok{(peso }\SpecialCharTok{\textasciitilde{}}\NormalTok{ dieta }\SpecialCharTok{+}\NormalTok{ genero, }\AttributeTok{data =}\NormalTok{ datos)}
\FunctionTok{summary}\NormalTok{(modelo)}

\CommentTok{\# SIN bloques (PEOR):}
\NormalTok{modelo\_malo }\OtherTok{\textless{}{-}} \FunctionTok{aov}\NormalTok{(peso }\SpecialCharTok{\textasciitilde{}}\NormalTok{ dieta, }\AttributeTok{data =}\NormalTok{ datos)}
\FunctionTok{summary}\NormalTok{(modelo\_malo)}
\end{Highlighting}
\end{Shaded}

\subsection{¿Quiero explorar muchos factores
eficientemente?}\label{quiero-explorar-muchos-factores-eficientemente}

\textbf{Ejemplo pregunta}: \emph{``Tienes 7 factores, saber cuáles son
importantes''}

\begin{Shaded}
\begin{Highlighting}[]
\FunctionTok{library}\NormalTok{(FrF2)}

\CommentTok{\# 7 factores, solo 8 experimentos (2\^{}(7{-}4))}
\NormalTok{design }\OtherTok{\textless{}{-}} \FunctionTok{FrF2}\NormalTok{(}\AttributeTok{nruns =} \DecValTok{8}\NormalTok{, }\AttributeTok{nfactors =} \DecValTok{7}\NormalTok{, }\AttributeTok{randomize =} \ConstantTok{FALSE}\NormalTok{)}

\CommentTok{\# Añadir respuesta}
\NormalTok{design }\OtherTok{\textless{}{-}} \FunctionTok{add.response}\NormalTok{(design, respuesta)}

\CommentTok{\# Analizar}
\NormalTok{modelo }\OtherTok{\textless{}{-}} \FunctionTok{lm}\NormalTok{(design, }\AttributeTok{degree =} \DecValTok{1}\NormalTok{)}
\FunctionTok{summary}\NormalTok{(modelo)}
\end{Highlighting}
\end{Shaded}

\newpage

\section{TABLA RESUMEN RÁPIDA}\label{tabla-resumen-ruxe1pida}

\begin{longtable}[]{@{}ll@{}}
\toprule\noalign{}
Pregunta del Examen & Test/Método \\
\midrule\noalign{}
\endhead
\bottomrule\noalign{}
\endlastfoot
¿Los hombres pesan más que las mujeres? & \textbf{t-test} \\
Hay outliers extremos + comparar grupos & \textbf{Wilcoxon} \\
Compara el peso entre 4 países & \textbf{ANOVA} \\
Compara 3 estados con outliers & \textbf{Kruskal-Wallis} \\
¿La variabilidad es igual en ambos grupos? & \textbf{F-test} \\
¿Ser del Barça está asociado con ser catalán? & \textbf{Fisher/Chi²} \\
¿Jugar en casa afecta a victorias? (2×2) & \textbf{Fisher/Chi²} \\
¿El peso del coche afecta al consumo? & \textbf{lm()} \\
¿Edad, BMI y fumar predicen gastos? & \textbf{lm()} múltiple \\
¿BMI predice eventos cardíacos (sí/no)? & \textbf{glm(binomial)} \\
¿BMI predice número de visitas (0,1,2,\ldots)? &
\textbf{glm(poisson)} \\
Visualiza 60 sensores en 2D & \textbf{PCA} \\
Muchas variables, cuáles son significativas & \textbf{FDR} \\
¿Los datos son normales? & \textbf{Shapiro-Wilk} \\
Comparar 3 dietas controlando por género & \textbf{Blocking} \\
¿Cuántos pacientes necesito? & \textbf{pwr.t.test()} \\
\end{longtable}

\section{CHECKLIST RÁPIDO}\label{checklist-ruxe1pido}

\textbf{PASO 1: ¿Qué me están preguntando?}

\begin{itemize}
\tightlist
\item[$\square$]
  Comparar grupos → Tests de hipótesis
\item[$\square$]
  Modelar relación → Regresión
\item[$\square$]
  Visualizar → PCA
\item[$\square$]
  Diseñar experimento → Power analysis
\end{itemize}

\textbf{PASO 2: Si es COMPARAR:}

\begin{itemize}
\tightlist
\item[$\square$]
  ¿Qué tipo? Continua/Categórica/Conteo
\item[$\square$]
  ¿Cuántos grupos? 1/2/3+
\item[$\square$]
  ¿Normal? → Shapiro
\item[$\square$]
  ¿Outliers? → Boxplot
\item[$\square$]
  ¿Independientes o relacionados?
\end{itemize}

\textbf{PASO 3: Si es MODELAR:}

\begin{itemize}
\tightlist
\item[$\square$]
  ¿Y es continua/binaria/conteo?
\item[$\square$]
  ¿Cuántos X?
\item[$\square$]
  Verificar supuestos después
\end{itemize}

\textbf{PASO 4: Si son MUCHAS variables:}

\begin{itemize}
\tightlist
\item[$\square$]
  Aplicar corrección FDR
\end{itemize}

\textbf{PASO 5: Si es VISUALIZAR:}

\begin{itemize}
\tightlist
\item[$\square$]
  PCA y colorear por grupos
\end{itemize}

\section{ERRORES COMUNES A EVITAR}\label{errores-comunes-a-evitar}

❌ Comparar dos grupos con outliers usando t-test\\
→ ✅ Usa Wilcoxon

❌ Usar ANOVA sin verificar normalidad\\
→ ✅ Primero Shapiro-Wilk, si no normal usa Kruskal-Wallis

❌ Modelar probabilidad de enfermedad con lm()\\
→ ✅ Usa glm(family=binomial)

❌ Probar 100 variables sin corrección múltiple\\
→ ✅ Aplica p.adjust(method=``BH'')

❌ Usar Chi² con frecuencias de 3\\
→ ✅ Usa Fisher para muestras pequeñas

❌ Hacer regresión sin graficar datos primero\\
→ ✅ Siempre plot(Y \textasciitilde{} X) primero

❌ PCA sin normalizar variables de diferentes escalas\\
→ ✅ Usa scale=TRUE en prcomp()

❌ Comparar dietas ignorando que género afecta\\
→ ✅ Usa diseño de bloques

\end{document}
